\documentclass[dvipdfmx, uplatex]{jsarticle}
\usepackage[dvipdfmx]{graphicx}
\usepackage{amsmath}
\usepackage{amssymb}
\usepackage{amsfonts}
\usepackage{amsthm}
\usepackage{mathrsfs}
\usepackage{bm}
\usepackage{braket}
\usepackage{xr}
\externaldocument{exercise_generating_functional}

\usepackage{silence}
\WarningFilter*{hyperref}{Token not allowed in a PDF string}
\usepackage[dvipdfmx]{hyperref}
\usepackage{pxjahyper}
\allowdisplaybreaks[4]

\usepackage[left=25truemm,top=20truemm,bottom=20truemm,right=25truemm]{geometry}

\theoremstyle{definition}
\newtheorem*{sol}{解答}

\title{解答：source 付き調和振動子の生成汎関数\\
\large ―― \texttt{exercise\_generating\_functional.tex} の各問の解答 ――}
\author{原典：戸板 真太郎『Lie 代数から学ぶ量子力学と場の量子論』\\
\texttt{LieQmQft}, \S「Generating Functional と Source Term」}
\date{}

\begin{document}
\maketitle

本ファイルは演習ノート \texttt{exercise\_generating\_functional.tex} の各問の完全な解答である。
記号・式番号は，特記なき限り問題ファイル \texttt{exercise\_generating\_functional.tex} のものを踏襲する
（\texttt{xr} パッケージにより，本文中の \eqref{eq:Wdef} 等は問題ファイルの式番号を指す）。
計算はすべて原典 \texttt{LieQmQft} の該当節を再構成したものであり，
個々のステップは Fresnel（Gauss）積分と三角関数の積和・半角公式に尽きる。

\tableofcontents

%======================================================================
\section{問題 1：被積分関数を Gauss 型に整理する}
%======================================================================

\begin{sol}
\par\smallskip\noindent\textbf{波動関数の積。}
\eqref{eq:vacuum} より $\phi_0^*(t_i+t,q_f)\,\phi_0(t_i,q_i)$ の位相は
\[
    +\frac{i(t_i+t)\omega}{2}-\frac{it_i\omega}{2}=+\frac{it\omega}{2}
\]
にまとまり，
\[
    \phi_0^*(t_i+t,q_f)\,\phi_0(t_i,q_i)
    =\left(\frac{m\omega}{\pi\hbar}\right)^{1/2}
    \exp\!\left(-\frac{m\omega}{2\hbar}(q_i^2+q_f^2)+\frac{it\omega}{2}\right).
\]
これと kernel \eqref{eq:kernel} を \eqref{eq:Wdef} に入れると，全体の前因子は
$\dfrac{m\omega}{\pi\hbar}(2i\sin\omega t)^{-1/2}$，位相は $e^{it\omega/2}$ となる。
あとは指数部 $-\dfrac{m\omega}{2\hbar}(q_i^2+q_f^2)+\dfrac{i}{\hbar}S_{cl}[J]$ を
$q_i,q_f$ について整理すればよい。

\par\smallskip\noindent\textbf{$\tfrac{i}{\hbar}S_{cl}$ の整理。}
\eqref{eq:Scl} に \eqref{eq:qcldot} の $\dot q_{cl}(t_i),\dot q_{cl}(t_f)$ と
\eqref{eq:qcl} の $q_{cl}$ を代入すると，$S_{cl}[J]$ は
$q_f^2,\,q_i^2,\,q_iq_f$（境界項由来），$q_f,\,q_i$（source と境界の結合），
および $J$ の 2 次項（$q$ 非依存）に分かれる。これらを
\[
    -\frac{m\omega}{2\hbar}(q_i^2+q_f^2)+\frac{i}{\hbar}S_{cl}[J]
    = -A(q_i^2+q_f^2)-2Bq_iq_f-2Cq_f-2Dq_i+E
\]
と書いたときの係数が以下である。

\par\smallskip\noindent\textbf{$A$（$q_i^2,q_f^2$ の係数）。}
波動関数由来の $\dfrac{m\omega}{2\hbar}$ と，$S_{cl}$ の境界項
$\dfrac12 m\,\dfrac{\omega\cos\omega t}{\sin\omega t}\,q^2$ に対応する
$-\dfrac{i}{\hbar}\dfrac{m\omega\cos\omega t}{2\sin\omega t}$ の和：
\begin{align}
    A
    =\frac{m\omega}{2\hbar}-\frac{i}{\hbar}\frac{m\omega\cos\omega t}{2\sin\omega t}
    =\frac{m\omega}{2\hbar}\left[1-i\frac{\cos\omega t}{\sin\omega t}\right]
    =\frac{m\omega}{2\hbar}\frac{-ie^{i\omega t}}{\sin\omega t}.
\end{align}
最後の等号は $1-i\cot\omega t=\dfrac{\sin\omega t-i\cos\omega t}{\sin\omega t}
=\dfrac{-i(\cos\omega t+i\sin\omega t)}{\sin\omega t}=\dfrac{-ie^{i\omega t}}{\sin\omega t}$。

\par\smallskip\noindent\textbf{$B$（$q_iq_f$ の係数の半分）。}
$S_{cl}$ の交差境界項 $-m\dfrac{\omega}{\sin\omega t}q_iq_f$ より
\begin{align}
    B
    =-\frac12\frac{i}{\hbar}\frac{-m\omega}{\sin\omega t}
    =\frac{m\omega}{2\hbar}\frac{i}{\sin\omega t}.
\end{align}

\par\smallskip\noindent\textbf{$C$（$q_f$ の係数の半分）。}
$q_f$ に線形に効くのは，$S_{cl}$ の境界項のうち source を含む部分と，
$\dfrac{q_f}{2}\int\cos\omega(t_i+t-\tau)J$ の項である。これらをまとめると
\begin{align}
    C
    &=-\frac12\frac{i}{\hbar}\frac12\int_{t_i}^{t_i+t}\!d\tau
    \left\{-\frac{m\omega}{\sin\omega t}\frac{\sin\omega(t_i+t-\tau)}{m\omega}\cos\omega t
    +\cos\omega(t_i+t-\tau)+\frac{\sin\omega(\tau-t_i)}{\sin\omega t}\right\}J(\tau)
    \notag\\
    &=-\frac{i}{4\hbar}\int_{t_i}^{t_i+t}\!d\tau\,
    \frac{-\sin\omega(t_i+t-\tau)\cos\omega t+\sin\omega t\cos\omega(t_i+t-\tau)+\sin\omega(\tau-t_i)}{\sin\omega t}\,J(\tau).
\end{align}
分子第 1・2 項は積和公式 $\sin\beta\cos\alpha-\cos\beta\sin\alpha=\sin(\beta-\alpha)$
（$\alpha=\omega t,\ \beta=\omega(t_i+t-\tau)$ に対し符号を整理）により
$\sin\omega(\tau-t_i)$ にまとまり，分子は $2\sin\omega(\tau-t_i)$ となる：
\begin{align}
    C
    =\frac{1}{2i\hbar}\frac{1}{\sin\omega t}\int_{t_i}^{t_i+t}\!d\tau\,\sin\omega(\tau-t_i)\,J(\tau)
    =\frac{1}{2i\hbar}\frac{\Im[e^{-i\omega t_i}F]}{\sin\omega t}.
\end{align}
ここで $\sin\omega(\tau-t_i)=\dfrac{e^{i\omega(\tau-t_i)}-e^{-i\omega(\tau-t_i)}}{2i}$ と
$F:=\int_{t_i}^{t_i+t}e^{i\omega\tau}J(\tau)\,d\tau$ を使い，
$\displaystyle\int\sin\omega(\tau-t_i)J=\frac{e^{-i\omega t_i}F-e^{i\omega t_i}F^*}{2i}=\Im[e^{-i\omega t_i}F]$ とした。

\par\smallskip\noindent\textbf{$D$（$q_i$ の係数の半分）。}
同様に，$q_i$ に線形に効く項をまとめると
\begin{align}
    D
    &=-\frac12\frac{i}{\hbar}\frac12\int_{t_i}^{t_i+t}\!d\tau
    \left\{\frac{m\omega}{\sin\omega t}\frac{\sin\omega(t_i+t-\tau)}{m\omega}
    +\frac{\sin\omega(t_i+t-\tau)}{\sin\omega t}\right\}J(\tau)
    \notag\\
    &=\frac{1}{2i\hbar}\int_{t_i}^{t_i+t}\!d\tau\,\frac{\sin\omega(t_i+t-\tau)}{\sin\omega t}\,J(\tau)
    =\frac{1}{2i\hbar}\frac{\Im[e^{i\omega(t_i+t)}F^*]}{\sin\omega t}.
\end{align}

\par\smallskip\noindent\textbf{$E$（$J$ の 2 次項）。}
$q_{cl}$ の source 部分どうしの積から来る $q$ 非依存項。\eqref{eq:Scl} の source 項
$\dfrac12\int q_{cl}J$ のうち $q_{cl}$ の source 部分を代入すると
\begin{align}
    E
    &=\frac{i}{2\hbar m\omega\sin\omega t}
    \int_{t_i<\tau'<\tau<t_i+t}\!\!\!d\tau\,d\tau'\,J(\tau)J(\tau')\,
    \big[\,\sin\omega(\tau-\tau')\sin\omega t
    \notag\\
    &\qquad
    -\sin\omega(\tau-t_i)\sin\omega(t_i+t-\tau')
    -\sin\omega(\tau'-t_i)\sin\omega(t_i+t-\tau)\,\big].
\label{eq:Eraw}
\end{align}
各 $\sin\cdot\sin$ を積和公式 $\sin a\sin b=\tfrac12[\cos(a-b)-\cos(a+b)]$ で $\cos$ に直すと
大半が打ち消し，
\begin{align}
    E
    =\frac{1}{i\hbar m\omega\sin\omega t}
    \int_{t_i}^{t_i+t}\!\!d\tau\int_{t_i}^{t_i+t}\!\!d\tau'\,
    \theta(\tau-\tau')\,J(\tau)J(\tau')\,
    \sin\omega(t_i+t-\tau)\sin\omega(\tau'-t_i)
\end{align}
にまとまる。さらに $\sin$ を指数で書き $F,F^*,|F|^2$ を導入すると
\begin{align}
    E
    &=\frac{i}{4\hbar m\omega\sin\omega t}
    \bigg\{
    2i\sin\omega t\!\int\!\!\int\! d\tau d\tau'\,J(\tau)J(\tau')\theta(\tau-\tau')e^{i\omega(\tau'-\tau)}
    \notag\\
    &\qquad\qquad
    -\tfrac12\Big(e^{i\omega(2t_i+t)}{F^*}^2+e^{-i\omega(2t_i+t)}F^2\Big)
    +e^{-i\omega t}|F|^2
    \bigg\}
\label{eq:Efinal}
\end{align}
となる（境界項 $F^2,{F^*}^2,|F|^2$ が現れることに注意。これは後で前因子側の項と相殺する）。
\end{sol}

%======================================================================
\section{問題 2：2 つの Gauss 積分}
%======================================================================

\begin{sol}
\eqref{eq:WABCD} の $q_f$ 積分から行う。$\Re A=\dfrac{m\omega}{2\hbar}\ge0$ なので
Fresnel 積分 \eqref{eq:fresnel} が使え，$q_f$ について平方完成して
\begin{align}
    \int dq_f\,e^{-Aq_f^2-2(Bq_i+C)q_f}
    =\sqrt{\frac{\pi}{A}}\,\exp\!\left(\frac{(Bq_i+C)^2}{A}\right).
\end{align}
残る $q_i$ 積分は
\begin{align}
    \int dq_i\,\exp\!\left(-Aq_i^2-2Dq_i+\frac{(Bq_i+C)^2}{A}\right)
    =\int dq_i\,\exp\!\left(-\Big(A-\tfrac{B^2}{A}\Big)q_i^2
    -2\Big(D-\tfrac{BC}{A}\Big)q_i+\tfrac{C^2}{A}\right),
\end{align}
であり，$\Re\!\big(A-\tfrac{B^2}{A}\big)=\Re\!\big(\tfrac{A^2-B^2}{A}\big)\ge0$
（問題 3 の \eqref{eq:A2mB2} と $\Re A>0$ から $0<\omega t<\pi$ で成立）なので再び
\eqref{eq:fresnel} を用いて
\begin{align}
    =\sqrt{\frac{\pi A}{A^2-B^2}}\,
    \exp\!\left(\frac{C^2}{A}+\frac{(AD-BC)^2}{A(A^2-B^2)}\right).
\end{align}
2 つの平方根と前因子 $\dfrac{m\omega}{\pi\hbar}(2i\sin\omega t)^{-1/2}$ を掛け合わせると
\begin{align}
    \frac{m\omega}{\pi\hbar(2i\sin\omega t)^{1/2}}\sqrt{\frac{\pi}{A}}\sqrt{\frac{\pi A}{A^2-B^2}}
    =\frac{m\omega}{\hbar[\,2i\sin\omega t\,(A^2-B^2)\,]^{1/2}}.
\end{align}
指数部は $\dfrac{C^2}{A}+\dfrac{(AD-BC)^2}{A(A^2-B^2)}
=\dfrac{C^2(A^2-B^2)+A(AD-BC)^2/A\cdots}{\cdots}$ を整理して
\[
    \frac{C^2}{A}+\frac{(AD-BC)^2}{A(A^2-B^2)}
    =\frac{A(C^2+D^2)-2BCD}{A^2-B^2}
\]
（分子を展開し $A^2-B^2$ でくくると確認できる）。よって \eqref{eq:Wgauss} を得る。
\end{sol}

%======================================================================
\section{問題 3：前因子の簡約}
%======================================================================

\begin{sol}
$A,B$ の値から
\begin{align}
    A^2-B^2
    =\left(\frac{m\omega}{2\hbar\sin\omega t}\right)^2\big[(-ie^{i\omega t})^2-(i)^2\big]
    =\left(\frac{m\omega}{2\hbar\sin\omega t}\right)^2\big[1-e^{2i\omega t}\big].
\end{align}
ここで $1-e^{2i\omega t}=-e^{i\omega t}(e^{i\omega t}-e^{-i\omega t})=-e^{i\omega t}\cdot 2i\sin\omega t$ なので
\begin{align}
    A^2-B^2
    =\left(\frac{m\omega}{2\hbar\sin\omega t}\right)^2(-2ie^{i\omega t}\sin\omega t)
    =\left(\frac{m\omega}{\hbar}\right)^2\frac{-ie^{i\omega t}}{2\sin\omega t},
\end{align}
これが \eqref{eq:A2mB2}。したがって
\begin{align}
    2i\sin\omega t\,(A^2-B^2)
    =2i\sin\omega t\left(\frac{m\omega}{\hbar}\right)^2\frac{-ie^{i\omega t}}{2\sin\omega t}
    =\left(\frac{m\omega}{\hbar}\right)^2 e^{i\omega t},
\end{align}
ゆえに
\begin{align}
    \frac{m\omega}{\hbar[\,2i\sin\omega t\,(A^2-B^2)\,]^{1/2}}
    =\frac{m\omega}{\hbar\cdot(m\omega/\hbar)\,e^{i\omega t/2}}
    =e^{-i\omega t/2}.
\end{align}
これが \eqref{eq:prefactor} である。
\end{sol}

%======================================================================
\section{問題 4：指数部の簡約}
%======================================================================

\begin{sol}
\par\smallskip\noindent\textbf{$C^2+D^2$。}
$C,D$ を $F,F^*$ で表した式（問題 1）より
\begin{align}
    C^2+D^2
    =-\frac{1}{4\hbar^2\sin^2\omega t}
    \left\{
    \Big(\frac{e^{-i\omega t_i}F-e^{i\omega t_i}F^*}{2i}\Big)^2
    +\Big(\frac{e^{i\omega(t_i+t)}F^*-e^{-i\omega(t_i+t)}F}{2i}\Big)^2
    \right\}.
\end{align}
2 乗を展開し $|F|^2=FF^*$ を用いると
\begin{align}
    C^2+D^2
    =\frac{\big(1+e^{-2i\omega t}\big)e^{-2i\omega t_i}F^2
    +\big(1+e^{2i\omega t}\big)e^{2i\omega t_i}{F^*}^2-4|F|^2}{16\hbar^2\sin^2\omega t}.
\end{align}

\par\smallskip\noindent\textbf{$BCD$。}
\begin{align}
    BCD
    =\frac{m\omega}{32\hbar^3}\frac{i}{\sin^3\omega t}
    \big[e^{-i\omega t_i}F-e^{i\omega t_i}F^*\big]
    \big[e^{i\omega(t_i+t)}F^*-e^{-i\omega(t_i+t)}F\big]
\end{align}
を展開すると
\begin{align}
    BCD
    =\frac{m\omega}{32\hbar^3}\frac{i}{\sin^3\omega t}
    \big[(e^{i\omega t}+e^{-i\omega t})|F|^2-e^{-i\omega(2t_i+t)}F^2-e^{i\omega(2t_i+t)}{F^*}^2\big].
\end{align}

\par\smallskip\noindent\textbf{$\dfrac{A(C^2+D^2)-2BCD}{A^2-B^2}$。}
$A=\dfrac{m\omega}{2\hbar}\dfrac{-ie^{i\omega t}}{\sin\omega t}$ を掛けて $2BCD$ を引くと
共通因子が括れ，
\begin{align}
    A(C^2+D^2)-2BCD
    =\frac{m\omega}{32\hbar^3}\frac{-ie^{i\omega t}}{\sin^3\omega t}
    \Big\{(1-e^{-2i\omega t})e^{-2i\omega t_i}F^2
    +(e^{2i\omega t}-1)e^{2i\omega t_i}{F^*}^2-2(1-e^{-2i\omega t})|F|^2\Big\}.
\end{align}
\eqref{eq:A2mB2} で割ると $\dfrac{1}{A^2-B^2}=\dfrac{\hbar^2}{(m\omega)^2}\dfrac{2\sin\omega t}{-ie^{i\omega t}}$ なので
\begin{align}
    \frac{A(C^2+D^2)-2BCD}{A^2-B^2}
    &=\frac{1-e^{-2i\omega t}}{16\hbar m\omega\sin^2\omega t}
    \Big\{e^{-2i\omega t_i}F^2+e^{2i\omega(t_i+t)}{F^*}^2-2|F|^2\Big\}
    \notag\\
    &=\frac{2ie^{-i\omega t}}{16\hbar m\omega\sin\omega t}
    \Big\{e^{-2i\omega t_i}F^2+e^{2i\omega(t_i+t)}{F^*}^2-2|F|^2\Big\},
\end{align}
ただし最後に $1-e^{-2i\omega t}=e^{-i\omega t}(e^{i\omega t}-e^{-i\omega t})=2ie^{-i\omega t}\sin\omega t$ を用いた。

\par\smallskip\noindent\textbf{$E$ との和：相殺。}
\eqref{eq:Efinal} の $E$ をこれに加えると，$F^2,{F^*}^2,|F|^2$ の境界項が
\[
    \underbrace{-2e^{i\omega(2t_i+t)}{F^*}^2-2e^{-i\omega(2t_i+t)}F^2+4e^{-i\omega t}|F|^2}_{E\text{ 側}}
    \;+\;
    \underbrace{2e^{-i\omega(2t_i+t)}F^2+2e^{2i\omega(t_i+t)}{F^*}^2-4e^{-i\omega t}|F|^2}_{\text{右項側}}
    =0
\]
のように完全に打ち消す（$2t_i+t$ の符号と位相が一致することを確認せよ）。
残るのは $E$ の 2 重積分項のみで，
\begin{align}
    E+\frac{A(C^2+D^2)-2BCD}{A^2-B^2}
    =-\frac{1}{2\hbar m\omega}
    \int_{t_i}^{t_i+t}\!\!d\tau\int_{t_i}^{t_i+t}\!\!d\tau'\,
    J(\tau)J(\tau')\,\theta(\tau-\tau')\,e^{i\omega(\tau'-\tau)}.
\end{align}
これが \eqref{eq:exponent} である。
\end{sol}

%======================================================================
\section{問題 5：最終形と Feynman 伝播関数}
%======================================================================

\begin{sol}
\par\smallskip\noindent\textbf{最終形。}
\eqref{eq:Wgauss} に問題 3・4 の結果を代入すると，前因子 $e^{-i\omega t/2}$ と位相 $e^{it\omega/2}$ が
相殺し（$e^{-i\omega t/2}\cdot e^{i\omega t/2}=1$），
\begin{align}
    W_{00}[J](t)
    =\exp\!\left(
    -\frac{1}{2\hbar m\omega}\int\!\!\int d\tau d\tau'\,JJ\,\theta(\tau-\tau')e^{i\omega(\tau'-\tau)}
    \right).
\end{align}
被積分関数を対称化する。積分は $\tau,\tau'$ の対称な領域でとられているので，
$\theta(\tau-\tau')e^{i\omega(\tau'-\tau)}$ の半分はそのまま，もう半分で $\tau\leftrightarrow\tau'$ と
変数を入れ替えると $\theta(\tau'-\tau)e^{i\omega(\tau-\tau')}$ が現れる：
\begin{align}
    -\frac{1}{2\hbar m\omega}\,\theta(\tau-\tau')e^{i\omega(\tau'-\tau)}
    \;\longrightarrow\;
    \frac{i}{4\hbar}\cdot\frac{i}{m\omega}
    \Big[\theta(\tau-\tau')e^{-i\omega(\tau-\tau')}+\theta(\tau'-\tau)e^{i\omega(\tau-\tau')}\Big].
\end{align}
よって \eqref{eq:DF} の $D_F$ を用いて \eqref{eq:Wfinal} を得る。

\par\smallskip\noindent\textbf{Fourier 表示。}
階段関数の表示
$\theta(s)=-\displaystyle\int\frac{d\omega'}{2\pi i}\frac{e^{-i\omega' s}}{\omega'+i\epsilon}$
を \eqref{eq:DF} の各項に入れる（$s=\tau-\tau'$）：
\begin{align}
    D_F
    =-\frac{i}{m\omega}\int\frac{d\omega'}{2\pi i}
    \left[\frac{e^{-i(\omega+\omega')(\tau-\tau')}}{\omega'+i\epsilon}
    +\frac{e^{i(\omega+\omega')(\tau-\tau')}}{\omega'+i\epsilon}\right].
\end{align}
第 1 項で $\omega'\to\omega'-\omega$，第 2 項で $\omega'\to-\omega'-\omega$ と平行移動すると分母が
$\omega'\mp\omega+i\epsilon$ となり，両者を通分して
\begin{align}
    D_F
    =\frac{i}{m\omega}\int\frac{d\omega'}{2\pi i}
    \frac{e^{-i\omega'(\tau-\tau')}}{\omega'^2-(\omega-i\epsilon)^2}
    \big[(-\omega'-\omega+i\epsilon)+(\omega'-\omega+i\epsilon)\big]
    =\frac{2}{m}\int\frac{d\omega'}{2\pi}
    \frac{e^{-i\omega'(\tau-\tau')}}{-\omega'^2+\omega^2-i\epsilon},
\end{align}
すなわち \eqref{eq:DFfourier}。

\par\smallskip\noindent\textbf{Green 関数。}
\eqref{eq:DFfourier} に $\dfrac{m}{2}(\partial_t^2+\omega^2)$ を作用させると，
$\partial_t^2\to-\omega'^2$ より
\begin{align}
    \frac{m}{2}\Big(\frac{d^2}{dt^2}+\omega^2\Big)D_F(t)
    =\int\frac{d\omega'}{2\pi}\frac{-\omega'^2+\omega^2}{-\omega'^2+\omega^2-i\epsilon}e^{-i\omega' t}
    \xrightarrow{\epsilon\to0}\int\frac{d\omega'}{2\pi}e^{-i\omega' t}
    =\delta(t),
\end{align}
これが \eqref{eq:green}。$D_F$ は確かに演算子 $\tfrac{m}{2}(\partial_t^2+\omega^2)$ の Green 関数である。
\end{sol}

%======================================================================
\section{問題 6：経路積分による系統的理解}
%======================================================================

\begin{sol}
\eqref{eq:Wdef} に kernel の経路積分表示 $K[J]=\int_{q(t_i)=q_i}^{q(t_f)=q_f}\mathcal{D}q\,e^{iS/\hbar}$ を入れ，
端点 $q_i,q_f$ の積分を波動関数と合わせると，端点を固定しない全経路積分
\begin{align}
    W_{FI}[J](t)
    =\int_{\forall q(t)}\mathcal{D}q\;\phi_F^*(t_i+t,q_f)\,e^{iS/\hbar}\,\phi_I(t_i,q_i)
    =\bra{\phi_F}\int_{\forall q(t)}\mathcal{D}q\,e^{iS/\hbar}\ket{\phi_I}
\end{align}
になる。source 付き調和振動子の作用は部分積分により
\begin{align}
    \frac{i}{\hbar}S
    =-\frac12\int d\tau\,q(\tau)\,\frac{i}{\hbar}m\Big(\frac{d^2}{d\tau^2}+\omega^2\Big)q(\tau)
    +\int d\tau\,q(\tau)\,\frac{i}{\hbar}J(\tau)
\end{align}
と書けるので，これは
$-\tfrac12 x_iA_{ij}x_j+x_iJ_i$（指数部）型の無限次元 Gauss 積分そのものである。
有限次元の公式
\[
    \int\prod_i\frac{dx_i}{\sqrt{2\pi}}\,e^{-\frac12 x_iA_{ij}x_j+x_iJ_i}
    =\frac{1}{\sqrt{\det A}}\,\exp\!\Big(\tfrac12 J_i(A^{-1})_{ij}J_j\Big)
\]
において，行列 $A\to\tfrac{i}{\hbar}m(\partial_\tau^2+\omega^2)$，source $J_i\to\tfrac{i}{\hbar}J(\tau)$，
逆行列 $A^{-1}\to$ Green 関数と読み替えれば，
\begin{align}
    W[J]
    &=(\text{規格化})\times
    \exp\!\left[\frac12\int d\tau d\tau'\,\frac{i}{\hbar}J(\tau')
    \Big\{\frac{i}{\hbar}m\big(\partial_\tau^2+\omega^2\big)\Big\}^{-1}\frac{i}{\hbar}J(\tau)\right]
    \notag\\
    &=\exp\!\left[\frac{i}{4\hbar}\int d\tau d\tau'\,J(\tau')
    \Big\{\frac{m}{2}\big(\partial_\tau^2+\omega^2\big)\Big\}^{-1}J(\tau)\right]
\end{align}
となる。$\big\{\tfrac{m}{2}(\partial_\tau^2+\omega^2)\big\}^{-1}=D_F$（問題 5 の \eqref{eq:green}）だから，
これは \eqref{eq:Wfinal} に一致する。規格化定数（$J$ 非依存部分）が
$W[0]=1$ を与えることは，問題 5 で前因子が完全に相殺したことに対応している。
\end{sol}

%======================================================================
\section{問題 7：分配関数と統計力学}
%======================================================================

\begin{sol}
波動関数 $\phi_0$ の寄与とは，$A$ に含まれる $\dfrac{m\omega}{2\hbar}$（および位相 $e^{it\omega/2}$）である。
これを除き $q_i=q_f=x$ とおくと，指数部は $-2Gx^2-2Hx$ となる。ここで
\begin{align}
    G
    &=A-\frac{m\omega}{2\hbar}+B
    =-\frac{i}{\hbar}\frac{m\omega\cos\omega t}{2\sin\omega t}
    -\frac12\frac{i}{\hbar}\frac{-m\omega}{\sin\omega t}
    =\frac{i}{\hbar}\frac{m\omega(1-\cos\omega t)}{2\sin\omega t}
    =\frac{i}{\hbar}\frac{m\omega\sin^2(\omega t/2)}{\sin\omega t},
    \\
    H
    &=C+D
    =\frac{1}{2i\hbar}\int_{t_i}^{t_i+t}\!d\tau\,
    \frac{\sin\omega(\tau-t_i)+\sin\omega(t_i+t-\tau)}{\sin\omega t}\,J(\tau).
\end{align}
$x$ について \eqref{eq:fresnel} を 1 回使い
\begin{align}
    Z[J](\tau)
    &=\left(\frac{m\omega}{2\pi i\hbar\sin\omega t}\right)^{1/2}\int dx\,e^{-2Gx^2-2Hx}
    =\left(\frac{m\omega}{2\pi i\hbar\sin\omega t}\right)^{1/2}\sqrt{\frac{\pi}{2G}}\,
    \exp\!\left(\frac{H^2}{2G}\right).
\end{align}
前因子は $G$ を代入して
\begin{align}
    \left(\frac{m\omega}{2\pi i\hbar\sin\omega t}\cdot
    \frac{\pi}{2\cdot\tfrac{i}{\hbar}\tfrac{m\omega\sin^2(\omega t/2)}{\sin\omega t}}\right)^{1/2}
    =\left(\frac{1}{4i^2\sin^2(\omega t/2)}\cdot\frac{1}{1}\right)^{1/2}
    =\frac{1}{2i\sin(\omega t/2)}
    =\frac{1}{e^{i\omega t/2}-e^{-i\omega t/2}},
\end{align}
よって \eqref{eq:Z}。$J=0$ では $H=0$ なので $\exp(H^2/2G)=1$，
虚時間 $\tau=it=\hbar\beta$ で $e^{i\omega t/2}=e^{\hbar\beta\omega/2}$ ゆえ
\begin{align}
    Z[J=0]
    =\frac{1}{e^{\beta\hbar\omega/2}-e^{-\beta\hbar\omega/2}}
    =\frac{1}{2\sinh(\beta\hbar\omega/2)}
    =\frac{e^{-\beta\hbar\omega/2}}{1-e^{-\beta\hbar\omega}}
    =\sum_{n=0}^{\infty}e^{-\beta\hbar\omega(n+1/2)},
\end{align}
これは統計力学で習う調和振動子（零点 energy $\tfrac12\hbar\omega$ を含む）の分配関数に他ならない。
\end{sol}

\end{document}
